\section{Fitting}\label{sec:fitting}

\subsection{Variable projection}\label{sec:varpro}

In equation~\ref{eq:bilinear} the spectra $\Smat$ enter linearly. For any fixed
set of non-linear parameters (the lifetimes, $t_0$ and the IRF width) the
best-fit spectra are therefore available in closed form,

\begin{equation}
  \Smat^{\mathsf{T}} = \left(\Cmat^{\mathsf{T}}\Cmat\right)^{-1}\Cmat^{\mathsf{T}}\Dmat ,
  \label{eq:varpro}
\end{equation}

and only the non-linear parameters need to go to the optimiser. This is
\emph{variable projection}: for a 500-pixel, three-component fit it reduces the
parameter vector from about 1500 entries to three, and it conditions the
problem far better than fitting amplitudes and lifetimes jointly. Amplitudes
must not be placed in the optimiser's parameter vector.

In practice equation~\ref{eq:varpro} is evaluated through a least-squares
solver rather than by forming the normal equations, which squares the condition
number of \Cmat.

\subsection{Engines}

\begin{table}[h]
  \centering
  \begin{tabular}{lll}
    \toprule
    Engine & Data & Output \\
    \midrule
    Single-trace fit & one kinetic trace       & lifetimes, amplitudes \\
    Global analysis  & full $\Dmat$, free \Kmat & lifetimes, DAS or EAS \\
    Target analysis  & full $\Dmat$, fixed connectivity & rates, SAS \\
    LDA              & full $\Dmat$, dense grid & lifetime density map \\
    \bottomrule
  \end{tabular}
  \caption{The four fitting engines. Global, target and single-trace share the
    variable-projection driver; LDA is a separate regularised linear solve.}
  \label{tab:engines}
\end{table}

\begin{lstlisting}
# Global analysis
fit = pr.fit_global(data, n_components=3, model_type="Sequential")
print(fit.taus, fit.tau_err)
data.plot_species_spectra(*fit.as_species_args())

# Preview (no optimisation)
prev = pr.preview_global(data, taus=[1.0, 10.0, 100.0])

# Target analysis
fit = pr.fit_target(data, scheme="A_eq_B_to_C", taus=[2.0, 5.0, 20.0, 200.0])

# Single-trace
fit = pr.fit_kinetics(t, y, n_components=2)

# Lifetime density analysis
lda = pr.fit.lda.solve_lda(data, n_taus=100, alpha="auto", non_negative=True)
\end{lstlisting}

Every fit routine is callable headlessly: a fit must be reachable and testable
from a plain script with no \code{QApplication}. The GUI is a caller, never a
prerequisite. Where a long fit puts up a progress dialog, the optimiser reports
through a plain callback and the GUI adapts it; the engine itself does not
touch widgets, and guards on \code{QApplication.instance()} being \code{None}.

\paragraph{Preview mode.} \code{pr.preview\_global} evaluates a model at
the given lifetimes \emph{without} optimising. The spectra still follow from
the data (they are the least-squares solution for the given \Cmat), so a
preview shows what those lifetimes actually imply and what the residual would
be. Nothing is fitted, so there is no convergence report and no uncertainties
--- and the result says so rather than presenting a guess as an answer.

\paragraph{Component ordering.} Variable projection sees only the column
space of \Cmat, so permuting the lifetimes fits equally well and the
optimiser's ordering depends on the initial guess. For parallel and sequential
models the components are sorted by lifetime on output, the spectra permuted to
match, and the fact recorded on the result (\code{sorted\_by\_lifetime}). A
target model is left exactly as fitted: its compartments are defined by the
scheme, and reordering them would misrepresent it.

\subsection{The coherent artefact}\label{sec:artefact}

Around time zero a transient carries signal no kinetic model describes:
cross-phase modulation and two-photon absorption in the solvent and the cell
windows. Left alone, the fit absorbs it into the shortest lifetime, which is
then not a lifetime at all.

Ticking \emph{Fit coherent artefact} (beside the IRF table, or
\code{coherent\_artifact=True} on any engine) appends three columns to the
design matrix: the instrument response and its first two derivatives,
\[
  g(t),\qquad g'(t),\qquad g''(t),
\]
each scaled to unit maximum so the second derivative --- larger than the
Gaussian by $1/\sigma^{2}$ --- does not dominate the conditioning. Their
amplitudes are solved per probe pixel exactly like a species spectrum, so the
artefact gets spectra of its own instead of distorting $\tau_1$. The
description follows Kovalenko \emph{et al.}, Phys.\ Rev.\ A \textbf{59}, 2369
(1999), which the log cites when the basis is first built.

The artefact columns are \emph{not} species: the result records how many there
are (\code{n\_artifact}), \code{as\_species\_args()} hands over only the kinetic
ones --- a column with no lifetime must not be labelled with someone else's ---
and \code{artefact\_spectra} exposes the rest. They are worth a look: an
artefact spectrum with amplitude far from time zero is a sign it is absorbing
real kinetics. The option needs a Gaussian IRF, since the basis is built from
it; without one the checkbox is disabled and the engines refuse rather than
inventing a width.

\subsection{Constrained amplitudes and sign masks}

When physical constraints are imposed on the species spectra --- non-negative
absorption for real excited states, non-positive for a ground-state bleach
contribution --- variable projection solves a bounded linear least-squares
problem per probe pixel rather than an unconstrained linear solve,

\begin{equation}
  \min_{\mathbf{S}_p}\,
    \lVert \mathbf{C}_w\,\mathbf{S}_p - \mathbf{d}_{w,p}\rVert^2
  \quad\text{subject to}\quad
    \mathbf{l} \le \mathbf{S}_p \le \mathbf{u}.
\end{equation}

A \code{sign\_mask} array of length $N_c$ controls this: $+1$ enforces
$S_p \ge 0$, $-1$ enforces $S_p \le 0$, and $0$ is unconstrained. When all
entries are zero the problem collapses to the standard unconstrained solve.
Simple non-negative least squares uses \code{scipy.optimise.nnls}; a mixed
mask uses \code{scipy.optimise.lsq\_linear}. The constrained solver is in
\code{pyrate_ta.fit.constrained}.

\subsection{Ground-state bleach and absolute species spectra}\label{sec:groundstate}

Transient absorption and vibrational data are difference spectra:
$\mathbf{S}_\text{diff} = \mathbf{S}_\text{excited} - f\,\mathbf{GS}$.
Adding back a scaled ground-state spectrum converts them to absolute absorption
spectra:
\begin{equation}
  \mathbf{S}_\text{abs}[:,\,i]
    = \mathbf{S}_\text{diff}[:,\,i] + f_i\cdot a\cdot \mathbf{GS},
\end{equation}
where $a$ is the ground-state scale factor and $f_i$ is the fraction of
ground state bleached by species $i$.

\code{pyrate_ta.fit.groundstate.compute\_absolute\_spectra} performs this
conversion. The maximum physically allowable scale factor $a_\text{max}$ ---
the largest $a$ before any pixel in any absolute spectrum turns negative --- is
computed first, so the automatic choice (\code{gs\_scale="auto"}) selects the
tightest non-negative solution. The result stores \code{S\_abs},
\code{gs\_scale} and \code{a\_max} for plotting and reporting.

\subsection{Lifetime density analysis}\label{sec:lda}

Instead of fitting $N$ discrete lifetimes, LDA solves for a continuous
distribution of amplitudes over a dense, logarithmically spaced grid of fixed
lifetimes,

\begin{equation}
  \min_{\Smat^{\mathsf{T}}}
  \left\lVert \mathbf{W} \odot \left(\Dmat - \Cmat(\boldsymbol{\tau}_\text{grid})\,
  \Smat^{\mathsf{T}}\right)\right\rVert^2
  + \alpha^2 \left\lVert \mathbf{L}\,\Smat^{\mathsf{T}}\right\rVert^2,
  \label{eq:lda}
\end{equation}

where $\mathbf{L}$ is a regularisation operator and $\alpha$ controls the
balance between fit quality and smoothness.

\paragraph{Regularisation operators.} Three are available: the identity
(ridge, $\mathbf{L} = \mathbf{I}$), a first-difference operator ($\mathbf{L} =
\mathbf{D}_1$, penalising slope), and a second-difference operator ($\mathbf{L}
= \mathbf{D}_2$, penalising curvature; the default). The augmented system
$[\Cmat;\,\alpha\mathbf{L}]\,\Smat^{\mathsf{T}} = [\Dmat;\,\mathbf{0}]$ is
solved in one least-squares call.

\paragraph{Automatic alpha selection.} When \code{alpha="auto"} three
methods are available:
\begin{description}
  \item[\code{lcurve}] The corner of the log-residual--log-norm curve ---
    maximum curvature by the Menger formula (Hansen 1992). This is the default.
  \item[\code{gcv}] Generalised cross-validation (Golub et al., 1979).
  \item[\code{morozov}] Morozov discrepancy principle: match the residual norm
    to the estimated noise level, estimated from the pre-time-zero region.
\end{description}

\paragraph{Additional options.}
\begin{description}
  \item[\code{non\_negative}] Enforce $S \ge 0$ per probe pixel via NNLS
    after the alpha scan. Physically motivated for some spectral regions.
  \item[\code{svd\_components}] Pre-filter the data matrix with the top $K$
    SVD components before constructing \Cmat, suppressing noise-dominated
    singular values.
  \item[\code{n\_bootstraps}] Monte Carlo resampling of the residuals:
    refits the LDA \code{n\_bootstraps} times and records the standard
    deviation of the integrated amplitude $A(\tau) = \sum_p |S_{p,i}|$ as a
    confidence estimate.
  \item[\code{find\_peaks}] Automatic peak centroid detection from
    $A(\tau)$, using \code{scipy.signal.find\_peaks} with a 5\% prominence
    threshold.
  \item[\code{coherent\_artifact}] Appends the same IRF-derivative basis as
    the discrete engines (section~\ref{sec:artefact}) to absorb time-zero
    contributions, leaving the regularised columns for the kinetics only.
\end{description}

The result is an \code{LDAResult} object (section~\ref{sec:results-lda})
carrying the density map, the alpha scan diagnostics, residuals, and optional
bootstrap errors. It exports the map to PyMORGAN's PDAT format via
\code{.to\_pdat(path)}.

\subsection{Watching a fit run}\label{sec:monitor}

A global fit blocks for seconds to minutes, and the cost alone says little
about what is going wrong. Setting \code{[gui] fit\_monitor\_every = N} opens a
window showing one bar per parameter, its current value above it and its name
below, refreshed every $N$ residual evaluations; $N = 0$ (the default) never
opens it. The scale is symlog by default
(\code{fit\_monitor\_scale}), which is the only one that shows a lifetime of
3000 and a time zero of $\approx 0$ on the same axis. Fixed parameters are drawn
greyed, so a bar that never moves is not mistaken for one that has converged.

The rate is a setting because each update costs a redraw \emph{and} an
event-loop pump: at $N = 1$ a fast fit spends more time drawing than fitting,
while $N = 5$--$10$ is enough to watch a lifetime run away to the edge of the
delay window. The drawing lives in \code{pyrate_ta.plot.monitor} and is Qt-free, so
the same bars --- and \code{plot\_parameter\_history}, the trajectory against
evaluation number --- can be produced from a script.

\subsection{Weighting and the fit statistic}\label{sec:weights}

If the dataset carries a per-point noise estimate --- a sibling \code{.pdatn}
file loaded alongside the \code{.pdat}, or the spread over single scans --- the
residuals can be weighted by it. PyMORGAN supplies the array through
\code{Dataset1D.noise\_array()}; PyRATE-TA never estimates noise itself.

With weights $w = 1/\sigma$, and $N$ the number of points actually fitted and
$p$ the number of free parameters, the statistic is a true reduced
chi-squared,

\begin{equation}
  \chi^{2}_{\nu} = \frac{1}{N-p}\sum_{n,k}
    \left(\frac{D_{nk} - \left(\Cmat\Smat^{\mathsf{T}}\right)_{nk}}{\sigma_{nk}}\right)^{2},
  \label{eq:chi2red}
\end{equation}

so that $\chi^{2}_{\nu} \approx 1$ means the residuals sit at the noise level,
$\gg 1$ that the model is too simple (or the noise underestimated), and
$\ll 1$ that the data are being overfitted (or the noise overestimated).
Without weights the reported statistic is the plain sum of squared residuals,
which carries no such interpretation and is not comparable between datasets.
The read-out in the GUI is therefore labelled \emph{SSR} or \emph{red.\
chi2} according to which one is in force, and never simply ``$\chi^2$''.

The parameter count $p$ includes the $N_p N_c$ linear amplitudes: variable
projection solves them in closed form, but it still spends those degrees of
freedom, so counting only the non-linear parameters would flatter the fit. The
value obtained with the non-linear-only denominator is reported alongside
(\code{FitStatistic.value\_nonlinear\_dof}), because the two differ by a large
factor and the choice should never be implicit.

Two practical points. Noise values are clipped from below at
\code{noise\_floor\_fraction} of the median, because a single near-zero
$\sigma$ would otherwise dominate the cost function entirely; and points with
non-finite or non-positive $\sigma$ are dropped from the fit along with their
data, with the count recorded on the result. Weights are diagonal: correlation
of the noise between pixels is ignored, so $\chi^{2}_{\nu}$ should not be read
as more rigorous than that assumption allows.

Requesting weights for a dataset that has no noise array raises. Falling back
silently to an unweighted fit while still calling the result a reduced
chi-squared would be a misreported statistic, which is worse than a refusal.

\subsection{The fixed/free mask}

A caller that knows only about the lifetimes --- the GUI reads its mask from the
lifetime table --- may pass a mask of ``n\_lifetimes'' entries while the
parameter vector is longer. Time zero and the IRF width are appended to that
vector \emph{only when they are being fitted}, so the entries that follow are
free by construction and the mask is padded rather than refused. A mask of any
\emph{other} length is still an error, naming the parameters expected: padding
must complete an incomplete answer, not reinterpret a wrong one.

A \code{fixed} mask may also be a sequence of parameter \emph{names}
(e.g.\ \code{["tau1", "t0"]}), in which case the driver matches by name and
the rest are free.

\subsection{Uncertainties}

One-sigma uncertainties are reported from the linearised covariance estimate at
the optimum. They are conditional on the model being correct, and they are not
confidence intervals for the choice of model: a three-component fit of
two-component data can return small uncertainties on three meaningless
lifetimes. Parameters held fixed receive no uncertainty and are flagged as
fixed, so that a fixed lifetime can never be mistaken for a well-determined
one.

\subsection{Convergence}\label{sec:convergence}

A non-converged fit raises, or is flagged on the result object; it is never
returned as though it had converged. The optimiser's own termination report
(the reason for stopping, the iteration count, the final cost) travels with the
result. Hitting \code{max\_iterations} is a non-convergence, not a result.

Parameters resting on a finite bound at termination are reported separately
(\code{FitReport.at\_bounds}); a bounded optimum is a constraint, not an
optimum, and downstream code should treat it as such.

\subsection{Results}\label{sec:results}

Three result classes share the common fields of \code{KineticFit}:

\begin{description}
  \item[\code{KineticFit}] Single-trace or multi-trace result: lifetimes,
    amplitudes, residuals, reproducibility payload.
  \item[\code{GlobalFit}] Adds nothing structurally (the spectra live in
    \code{S}), but is typed distinctly so a function can assert it received
    a global result rather than a single-trace one.
  \item[\code{TargetFit}] Adds \code{scheme\_key}, \code{scheme\_label},
    \code{scheme\_text}, the rate matrix \code{K} and \code{eigenvalues}.
    The eigenvalues are stored because they, not the rates, are what the data
    actually determine for a branched scheme.
\end{description}

Every result exposes \code{as\_species\_args()} (the adapter for PyMORGAN's
\code{plot\_species\_spectra}), \code{spectra\_kind} (DAS / EAS / SAS),
\code{artefact\_spectra} (the coherent-artefact columns, if any),
\code{summary()} and \code{log\_summary()}.

\subsubsection{LDA result}\label{sec:results-lda}

\code{LDAResult} carries the density map \code{S\_map} of shape
$(N_p \times M)$, the lifetime grid \code{tau\_grid} of length $M$, the chosen
$\alpha$, the L-curve diagnostics, optional bootstrap errors, and a list of
detected peaks. It also carries \code{citations()} and \code{format\_citations()},
which enumerate the references for whichever algorithm combination was used, and
\code{to\_pdat(path)}, which exports the 2D map to PyMORGAN's PDAT format for
further plotting.

\subsection{Reproducibility}\label{sec:repro}

A result object carries enough to reproduce itself:

\begin{itemize}
  \item the model identity (family, number of components, the \Kmat{} where applicable);
  \item the initial guesses and the bounds;
  \item which parameters were held fixed;
  \item the delay and probe ranges actually fitted, including any clipping applied for the coherent artefact;
  \item the solver's convergence report.
\end{itemize}

A lifetime reported without its uncertainty and its fixed/free flag is not a
result. Results serialise through \code{pyrate_ta.io}, which handles fit sessions
--- parameters, masks, results --- and not raw spectroscopy formats.

\subsection{Fit session files}\label{sec:sessions}

\code{pr.save\_fit(path, fit)} writes a \code{.prfit} file: a compressed
\code{.npz} archive holding the arrays (spectra, concentration profiles,
residuals, axes) and the reproducibility payload as a JSON string inside it.
Keeping them together means a saved fit cannot lose its provenance.

\code{pr.load\_fit(path)} reopens a session as a \code{LoadedFit}: the arrays
as attributes, the JSON payload as \code{.meta}. It is deliberately \emph{not}
a live \code{KineticFit} --- a reopened session has no dataset behind it and
cannot be refitted as it stands, and returning something that looks like a
fresh result would invite exactly that. Everything needed to set the fit up
again is in \code{.meta}. A target fit reopens with its \code{scheme\_text},
so the scheme is editable without re-entering it.

\begin{lstlisting}
path = pr.save_fit("scan_fit", fit)     # -> scan_fit.prfit
prev = pr.load_fit("scan_fit.prfit")
print(prev.summary())                    # same format as a live result
print(prev.taus)                         # numpy arrays as attributes
\end{lstlisting}

The format version is embedded in every file and checked on load: a file
written by a newer PyRATE-TA that the current build cannot interpret raises rather
than silently mis-reading it.

\subsection{Testing}

Synthetic-recovery tests are the backbone of the test suite: build data from
known parameters, fit it, and assert that the parameters come back within
tolerance. A fitting project without these is untested regardless of its line
coverage. The failure modes are tested too --- a non-converged fit must be
flagged rather than silently returned, and near-degenerate lifetimes must not
produce exploding amplitudes without a warning.
